\chapter{Building a Differential Diagnosis}

\label{ch:ch34_differential_dx}

\section{Building a Differential}
Start with the imaging pattern and clinical context, then rank diagnoses by probability, urgency, and actionability. Include common diagnoses, dangerous alternatives, and relevant mimics without producing an indiscriminate list.

\subsection{The Clinical Framework}
Before analysing the images, consider:
\begin{itemize}
    \item \textbf{Age} -- a lytic bone lesion in a 25-year-old differs from one in a 70-year-old.
    \item \textbf{Sex} -- autoimmune hepatitis is more common in women; haemochromatosis in men.
    \item \textbf{History} -- prior malignancy, immunosuppression, surgery, and travel history all reshape the differential.
    \item \textbf{Acuity} -- an acute presentation demands urgent diagnoses (hemorrhage, ischemia, perforation); a chronic presentation permits broader differentials.
\end{itemize}

\subsection{The Imaging-Based Approach}
Once clinical context is established, systematically evaluate:
\begin{enumerate}
    \item \textbf{Location} -- hepatic segment, lung lobe, brain lobe, vertebral level. Certain pathologies have predilections: oligodendroglioma for frontal lobe, neuroblastoma for adrenal, chordoma for clivus.
    \item \textbf{Morphology} -- solid, cystic, fat-containing, calcified, infiltrative. A fat-containing lesion in the kidney is angiomyolipoma until proven otherwise.
    \item \textbf{Enhancement pattern} -- rim enhancement suggests abscess or necrotic tumor; homogeneous enhancement suggests a benign tumor or vascular lesion.
    \item \textbf{Diffusion restriction} -- on brain MRI, restricted diffusion in a vascular territory is acute stroke; in a mass, it suggests high cellularity (lymphoma, high-grade glioma).
    \item \textbf{Vascularity} -- on CT/MRI, a hypervascular liver lesion in a young woman on oral contraceptives is a hepatic adenoma; in a cirrhotic patient, it is hepatocellular carcinoma.
\end{enumerate}

\begin{pearlbox}[Clinical Pearl: Most-likely-first vs. worst-first]
In low-risk patients, start with the most likely diagnosis and explain why less likely diagnoses are excluded. In high-risk patients (immunocompromised, known malignancy, acute deterioration), always consider the worst plausible diagnosis first and rule it out before settling on a benign explanation.
\end{pearlbox}

\section{Discriminating Features}
For each leading possibility, identify the feature that would increase or decrease its likelihood: distribution, tissue behavior, enhancement, diffusion, vascularity, time course, or associated findings. Recommend the next test only when it addresses a specific unresolved question.

\subsection{How to Narrow a Differential}

\begin{table}[h]
\centering
\caption{Imaging features and their discriminating power.}
\label{tab:discriminating_features}
\begin{tabular}{p{4cm} p{4.5cm} p{5cm}}
\toprule
\textbf{Feature} & \textbf{Favours Diagnosis A} & \textbf{Favours Diagnosis B} \\
\midrule
Calcification in a lung nodule & Granuloma (benign, old) & Malignancy (rare; consider carcinoid or mucinous adenocarcinoma) \\
DWI restriction in a brain lesion & Acute stroke (vascular territory) & Abscess (ring-enhancing with central restriction) \\
Fat in a renal lesion & Angiomyolipoma (benign) & Renal cell carcinoma (rarely contains fat) \\
Peripheral enhancement of a liver lesion & Hepatic abscess or metastasis & Hemangioma (nodular peripheral enhancement with centripetal fill-in) \\
T2 shine-through on MRI & Gliosis or vasogenic edema & Acute ischemia (true restriction on ADC map) \\
\bottomrule
\end{tabular}
\end{table}

\section{The Role of Clinical Correlation}
Imaging rarely provides a definitive diagnosis in isolation. A solitary pulmonary nodule in a 60-year-old smoker with haemoptysis demands a different approach than the same nodule in a 30-year-old non-smoker. Laboratory values, clinical examination, and response to treatment are integral to the diagnostic process.

\begin{itemize}
    \item \textbf{Laboratory correlation} -- elevated alpha-fetoprotein in a cirrhotic patient with a hypervascular liver lesion confirms HCC.
    \item \textbf{Clinical correlation} -- a solitary thin-walled cavity in a patient with haemoptysis is cavitary lung cancer until proven otherwise, even though a simple cyst is statistically more common.
    \item \textbf{Treatment response} -- administering antibiotics and repeating imaging in 6--8 weeks is a diagnostic test: resolution favours infection, persistence favours neoplasm.
\end{itemize}

\begin{pearlbox}[Clinical Pearl: The undiagnosed mass]
When a mass cannot be characterized by imaging alone, state the differential in order of probability and recommend the most specific next step (biopsy, PET-CT, or targeted follow-up). Avoid ``nonspecific mass'' as a final impression without a structured differential.
\end{pearlbox}

\section{Communication}
Use calibrated language such as likely, possible, indeterminate, or highly suspicious. State when pathology, laboratory correlation, clinical examination, or follow-up is needed. Avoid definitive language when the imaging evidence is insufficient.

\subsection{Structuring the Impression}
A well-constructed impression:
\begin{enumerate}
    \item Leads with the most clinically important finding.
    \item Provides a ranked differential (not just a list).
    \item States the confidence level (likely, possible, indeterminate).
    \item Recommends a specific next step when appropriate.
    \item Uses standardised language (BI-RADS, Lung-RADS, etc.) when applicable.
\end{enumerate}

\begin{pearlbox}[Impression:]
A useful differential diagnosis is short (3--5 entries), ranked by probability, informed by clinical context, and communicated with calibrated uncertainty. The goal is to help the referring clinician decide what to do next, not to list every entity in the textbook.
\end{pearlbox}

\clearpage
